\uuid{LkhI}
\titre{Système linéaire à paramètres/Méthode du Pivot de Gauss}
\auteur{hanani ; mzari}
\organisation{exo7}
\licence{CC BY-NC-SA 4.0}
\datecreate{2020-01-20}
\niveau{L1}
\module{Algèbre}
\chapitre{Déterminant, système linéaire}
\sousChapitre{Système linéaire, rang}
\difficulte{3}
\theme{Systèmes d'équations linéaires, Système linéaire à paramètres, Méthode du Pivot de Gauss}
\exo7qcmid{lille-2.systemes.013}

\contenu{
\begin{question}
\qclassification{200.03}
\qtags{motcle=Système linéaire à paramètres/Méthode du Pivot de Gauss}
On considère le système d'équations, d'inconnue $(x,y,z,t) \in \Rr^4$ et  de paramètres des réels  $a$ et $m$ :
$$(\mathtt{S})
\left\{\begin{array}{rcc}
x-z-t&=&0\\
-x+y+z&=&a\\
2x+y-z&=&m \\
x-mz-t&=&a \\
x+y+t&=&m.\end{array}\right.$$
Quelles sont les assertions vraies ?
\begin{answers}
\good{$(\mathtt{S}) \Leftrightarrow
\left\{\begin{array}{rcc}
x-z-t&=&0\\
y-t&=&a\\
z+3t&=&m-a \\
(1-m)z&=&a.\end{array}\right.$}
\bad{Si $m=1$ et $a=0$,  $(\mathtt{S})$ admet une unique solution.}
\good{Si $m \neq 1$ et $a$ un réel quelconque,  $(\mathtt{S})$ admet une unique solution.}
\bad{Si $m \neq 1$ et $a\neq 0$, $(\mathtt{S})$ admet une infinité de solutions.}
\end{answers}
\begin{explanations}
$$(\mathtt{S}) \Leftrightarrow
\left\{\begin{array}{rcc}
x-z-t&=&0\\
y-t&=&a\\
y+z+2t&=&m \\
(1-m)z&=&a \end{array}\right.
\Leftrightarrow
\left\{\begin{array}{rcc}
x-z-t&=&0\\
y-t&=&a\\
z+3t&=&m-a \\
(1-m)z&=&a.\end{array}\right.$$
\begin{enumerate}
\item[-]Si $m=1$ et $a\neq 0$, $(\mathtt{S})$ n'admet pas  de solution.
\item[-]Si $m=1$ et $a= 0$, $(\mathtt{S})$ admet une infinité de solutions.
\item[-]Si $m\neq1$, $(\mathtt{S})$ admet une unique solution.
\end{enumerate}
\end{explanations}
\end{question}
}
